% ===================================================================
%  TABEL Nr. 15 — Turg. --- Toiduained.
%  Traduction 2026 d'apres texto/io, controlee sur texto/fr et
%  texto/fi. Consigne : voir texto/et/10-tabel-01.tex.
%
%  LA POMME DE TERRE DE L'ALINEA 10 EST LA SEPTIEME REPONSE DE LA
%  SERIE, ET C'EST ELLE QUI DONNE A LA QUESTION SA FORME DEFINITIVE.
%
%    lituanien      « bulvė »     — un emprunt ;
%    tcheque        « brambory »  — le nom du pays d'ou elle vient ;
%    allemand       « Kartoffel » — le nom italien de la truffe ;
%    francais       « pomme de terre » — un compose fait chez soi ;
%    luxembourgeois « Gromper »   — un autre compose fait chez soi ;
%    romanche       « tartuffel » — le mot allemand ;
%    ESTONIEN       « kartul »    — le mot allemand.
%
%  DEUX FAMILLES DE LANGUES, DEUX EMPRUNTS A LA MEME SOURCE. Le
%  tableau 15 romanche avait conclu que, quand on emprunte, c'est la
%  route de la chose qui choisit le preteur ; l'estonien le confirme
%  d'un tout autre bout de l'Europe, et pour la meme raison : la
%  pomme de terre y est arrivee du monde allemand.
%
%  MAIS L'ESTONIEN AJOUTE CE QUE LE ROMANCHE NE POUVAIT PAS DONNER,
%  ET C'EST DECISIF. Les dialectes estoniens ont eu aussi
%  « maaõun », la pomme-de-terre, compose fait a la maison sur le
%  meme modele que le francais et le luxembourgeois. LES DEUX FORMES
%  ONT COEXISTE ET C'EST L'EMPRUNT QUI A GAGNE.
%
%  Pourquoi ? Parce que le tubercule n'est pas arrive en Estonie par
%  les jardins mais PAR LE MANOIR : ce sont les proprietaires
%  baltes-allemands qui, au XVIIIe siecle, ont distribue les plants
%  et impose la culture. C'est le meunier du tableau 7 une seconde
%  fois — QUI DISTRIBUE NOMME.
%
%  LA SERIE DE LA POMME DE TERRE SE FERME DONC AINSI : SON NOM DIT
%  QUI VOUS A DONNE LE PREMIER PLANT. La ou le paysan l'a recu d'un
%  autre paysan qui parlait sa langue, il l'a decrite — pomme de
%  terre, Gromper, maaõun. La ou il l'a recu d'un maitre qui parlait
%  une autre langue, il a garde le mot du maitre — kartul,
%  tartuffel.
% ===================================================================

%%K t15-apar-1 apar
\VUsaut{4.0mm}
\VUtitre{15.0pt}{60}{TABEL Nr. 15}
\VUsaut{3.0mm}

%%K t15-tit-1 sub
\VUpk{11.4pt}{40}{Turg. --- Toiduained.}
\VUsaut{3.0mm}

%%K t15-01-1 p
1. --- Enamik kaupmehi, kes tarnivad meile meie toiduks vajalikku kaupa, kogunevad
kaetud turu \VUgras{halli} \textsuperscript{(1)} alla või naabertänavatele.

%%K t15-02-1 p
2. --- \VUgras{Sealihunik} \textsuperscript{(2)}, kes müüb \VUgras{sinki}
\textsuperscript{(3)}, \VUgras{verivorsti} \textsuperscript{(4)}, \VUgras{vorstikest}
\textsuperscript{(5)}, \VUgras{rupsivorsti} \textsuperscript{(6)} ja \VUgras{pekki}
\textsuperscript{(7)}, seisab \VUgras{või- ja juustumüüja} \textsuperscript{(8)} kõrval,
kelle väljapanek on varustatud punase koorikuga \VUgras{hollandi juustudega}
\textsuperscript{(9)}, \VUgras{camembert'i juustudega} \textsuperscript{(10)},
\VUgras{gruyère'i ketastega} \textsuperscript{(11)} ja värskete \VUgras{võitükkidega}
\textsuperscript{(12)}.

%%K t15-03-1 p
3. --- Tema ees pakub \VUgras{juurviljamüüja} \textsuperscript{(13)} oma klientidele
ilusaid \VUgras{tomateid} \textsuperscript{(14)}, \VUgras{lillkapsaid}
\textsuperscript{(15)}, \VUgras{salatit} \textsuperscript{(16)}, \VUgras{peakapsaid}
\textsuperscript{(17)}, \VUgras{kõrvitsaid} \textsuperscript{(19)}, \VUgras{meloneid}
\textsuperscript{(18)} ja isegi \VUgras{seeni} \textsuperscript{(20)}. Kuid õllepoiss, kes
on peatanud oma \VUgras{veovankri} \textsuperscript{(21)}, mis on laaditud
\VUgras{õllevaatidega} \textsuperscript{(22)}, just tema väljapaneku ees, takistab tema
klientidel ligi pääseda. Üks aiapidaja toob talle korvi \VUgras{artišokke}
\textsuperscript{(23)}.

%%K t15-tit-2 sub
\VUsaut{4.0mm}
\VUpk{10.2pt}{40}{Müümine ja ostmine.}
\VUsaut{3.0mm}

%%K t15-04-1 p
4. --- Turul on elevus suur, sest on saabunud \VUgras{oksjonimüügi}
\textsuperscript{(24)} tund. \VUgras{Perenaised} \textsuperscript{(26)}, kes tulevad siia
oma ostude pärast, peatuvad möödudes \VUgras{rupsimüüja} \textsuperscript{(25)} poe ees,
et osta \VUgras{rupsi} või \VUgras{vasikapead}.

%%K t15-05-1 p
5. --- Priske \VUgras{kokk} \textsuperscript{(27)} tingib väga ilusa
\VUgras{tuunikala} pärast \VUgras{kalamüüja} \textsuperscript{(28)} juures, kes oma
suva järgi hindu tõstab. Ta sai suure koguse \VUgras{angerjaid, merikeeli, forelle} ja
\VUgras{karpkalu}. Tema \VUgras{tursk} ja tema \VUgras{rannakarbid} on väga värsked.

%%K t15-tit-3 sub
\VUsaut{4.0mm}
\VUpk{10.2pt}{40}{Odav ja kallis kaup.}
\VUsaut{3.0mm}

%%K t15-06-1 p
6. --- \VUgras{Linnumüüja} \textsuperscript{(29)}, kes on seatud tema kõrvale, on väga
kohusetundlik. Kui ta müüb \VUgras{kalkunit, hane, parti} või lihtsat
\VUgras{kanapoega}, nõuab ta ainult õiglast hinda.

%%K t15-07-1 p
7. --- Platsi nurgas, teisel korrusel, on lühikest aega tagasi sisse seatud
\VUgras{riidemüüja} \textsuperscript{(30)}, kelle juurest ostjannad on alati kindlad
leidma väga ilusat valikut \VUgras{laudlinu, salvrätikuid, põllesid} ja
\VUgras{käterätikuid}. Samal korrusel on \VUgras{noassepp} \textsuperscript{(31)} sisse
seadnud oma töökoja. Ta teritab halli müüjannade \VUgras{nuge} ja valmistab aednikele
\VUgras{sirpe} kõige kõvemast \VUgras{terasest}.

%%K t15-tit-4 sub
\VUsaut{4.0mm}
\VUpk{10.2pt}{40}{Vürtspood.}
\VUsaut{3.0mm}

%%K t15-08-1 p
8. --- Tema töökoja all on lühikest aega tagasi avatud suur toiduainete kauplus. See ei
ole enam endine lihtne ja tähtsusetu \VUgras{vürtspood} \textsuperscript{(32)}, mis müüs
ainult üldkasutatavat kaupa, \VUgras{teed} \textsuperscript{(33)}, \VUgras{šokolaadi}
\textsuperscript{(34)}, \VUgras{pea suhkrut} \textsuperscript{(37)} ja \VUgras{seepi}
\textsuperscript{(38)}. Seal müüakse igasugust kaupa, \VUgras{krõbedaid küpsiseid},
\VUgras{konserve} \textsuperscript{(35)}, \VUgras{likööre} \textsuperscript{(36)},
\VUgras{rummi, konjakit, kirsiviina, aniisilikööri} ja \VUgras{curaçaod}. Seal leiab
ulukiliha: \VUgras{jänest} \textsuperscript{(39)}, \VUgras{rästaid}
\textsuperscript{(40)}, \VUgras{nurmkanu} \textsuperscript{(41)}, \VUgras{vutte}
\textsuperscript{(42)}, \VUgras{metsparti} \textsuperscript{(43)} või \VUgras{faasanit}
\textsuperscript{(44)}, sama kergesti kui eksootilisi puuvilju, nagu \VUgras{banaanid}
\textsuperscript{(46)} ja \VUgras{ananassid}.

%%K t15-09-1 p
9. --- Vürtspoodniku väga teenistusvalmis \VUgras{poiss} \textsuperscript{(45)} on valmis
andma seda, mida soovitakse: sõstra- või küdoonia\VUgras{moosi} \textsuperscript{(47)},
\VUgras{rasva} \textsuperscript{(48)}, \VUgras{kurgikesi} \textsuperscript{(49)} või
soolatud \VUgras{sardiine} \textsuperscript{(50)}.

%%K t15-09-2 p
See on väga mugav \VUgras{klientide} \textsuperscript{(51)} jaoks, kes leiavad, kõige
tavalisema kauba kõrvalt, kõige haruldasemad esimesed viljad ja kõige maitsvamad
puuviljad: mahlased \VUgras{pirnid} \textsuperscript{(52)}, \VUgras{ploomid}
\textsuperscript{(53)}, \VUgras{viinamarjad} \textsuperscript{(54)}, \VUgras{virsikud}
\textsuperscript{(55)}, \VUgras{mandlid} \textsuperscript{(56)} ja \VUgras{pähklid}
\textsuperscript{(57)}.

%%K t15-tit-5 sub
\VUsaut{4.0mm}
\VUpk{10.2pt}{40}{Kliendid.}
\VUsaut{3.0mm}

%%K t15-10-1 p
10. --- \VUgras{Riis} \textsuperscript{(58)} ja \VUgras{läätsed} \textsuperscript{(59)} on
suitsutatud \VUgras{heeringate} \textsuperscript{(60)} kasti kõrval; \VUgras{kartul}
\textsuperscript{(61)} ja \VUgras{oad} \textsuperscript{(63)} on \VUgras{õunte}
\textsuperscript{(62)}, \VUgras{viigimarjade} \textsuperscript{(64)}, \VUgras{kuivatatud
ploomide} \textsuperscript{(65)} ja \VUgras{sarapuupähklite} \textsuperscript{(66)}
naabruses. Selles kaupluses leiab värskeid \VUgras{mune} \textsuperscript{(67)} ja
\VUgras{oliive} \textsuperscript{(68)}; seepärast täituvad \VUgras{korvid}
\textsuperscript{(70)} ja \VUgras{toiduvõrgud} \textsuperscript{(73)} väga kiiresti.

%%K t15-11-1 p
11. --- Selle suurepärase firma \VUgras{kohv} \textsuperscript{(72)} on omandanud
\VUgras{teenijate} \textsuperscript{(69)} ja \VUgras{kokkade} seas õiglase kuulsuse, sest
vana \VUgras{vürtspoodnik} \textsuperscript{(71)} ise röstib seda kõige hoolikamalt.

%%K t15-tit-6 sub
\VUsaut{4.0mm}
\VUpk{10.2pt}{40}{Vaatepildid.}
\VUsaut{3.0mm}

%%K t15-12-1 p
12. --- Kõik ei lähe turule end varustama. Maalilises liikumises tänavakesel, mis viib
halli, näeb kõige mitmekesisemaid inimesi. Lahke \VUgras{lihuniku poisi}
\textsuperscript{(75)} kõrval, kes lükkab enda ees \VUgras{käru} \textsuperscript{(74)},
on siin kaks puhkusel sõdurit, väga uhked oma säravat vormi näidates:
\VUgras{husaar} \textsuperscript{(76)} sinise \VUgras{dolmani} ja \VUgras{sulgedega
tšakoga}, ja \VUgras{zuaavide kapral} pungis pükstega ja peas \VUgras{fess}.

%%K t15-13-1 p
13. --- \VUgras{Pagar} \textsuperscript{(80)}, kes seisab \VUgras{pagariäri}
\textsuperscript{(78)} lävel, on äsja sõtkunud oma \VUgras{leiva} \textsuperscript{(79)}
taigna ja võtnud ahjust välja \VUgras{sarvesaiad} \textsuperscript{(81)},
\VUgras{saiakesed} \textsuperscript{(82)} ja \VUgras{ümarad leivad} \textsuperscript{(83)},
mis on vitriinis.

%%K t15-14-1 p
14. --- Pagariäri kohal elab osav \VUgras{pesuõmbleja} \textsuperscript{(84)}, kes annab
tööd mitmele \VUgras{tikkijannale} \textsuperscript{(85)}. Tema kõrval elab
\VUgras{pesija-triikija} \textsuperscript{(86)}, kes tärgeldab \VUgras{pesu}
\textsuperscript{(88)}, olles selle pesnud ja kuivatanud, ja kes triigib seda
\VUgras{triikrauaga} \textsuperscript{(87)}.

%%K t15-tit-7 sub
\VUsaut{4.0mm}
\VUpk{10.2pt}{40}{Liha, kalad ja juurviljad.}
\VUsaut{3.0mm}

%%K t15-15-1 p
15. --- \VUgras{Lihapoes} \textsuperscript{(89)}, mis asub esimesel korrusel, müüakse
igasugust \VUgras{liha}, \VUgras{lambaliha} \textsuperscript{(90)}, \VUgras{veiseliha}
\textsuperscript{(94)} või \VUgras{vasikaliha} \textsuperscript{(92)}, nagu
\VUgras{lambakintse, biifsteeke, praade} ja \VUgras{kotlette}. \VUgras{Lihuniku poiss}
\textsuperscript{(93)} lõikab kõik need lihad lahti ja paneb eraldi \VUgras{üdikondid}
\textsuperscript{(96)}. Lihapoe ukse juures on \VUgras{austrimüüja}
\textsuperscript{(97)}, kes müüb \VUgras{austreid} \textsuperscript{(98)}. Ta avab need,
kui seda soovitakse.

%%K t15-16-1 p
16. --- Kõik need kaupmehed maksavad patendi- või kohatasu; seepärast jälitab
\VUgras{politseinik} \textsuperscript{(100)} halastamatult vaeseid
\VUgras{juurviljarändmüüjaid} \textsuperscript{(99)}, kes neile konkurentsi teevad,
kandes oma kärudel ringi \VUgras{porgandeid, rediseid, naereid, porrulauke} või
\VUgras{sparglikimpe, värskeid herneid, spinatit, hapuoblikat, salatkressi} ja
\VUgras{peterselli}.

%%K t15-tit-8 sub
\VUsaut{4.0mm}
\VUpk{10.2pt}{40}{Ettevaatust politseiniku eest!}
\VUsaut{3.0mm}

%%K t15-17-1 p
17. --- Poiss, kes müüb \VUgras{küüslauku} \textsuperscript{(101)} ja \VUgras{sibulaid},
teab väga hästi, et ka tema ei tohi oma kaupa turu lähedal müüa. Ta põgeneb, joostes
riskiga ümber ajada \VUgras{krabide}, \VUgras{vähkide} ja \VUgras{krevettide} korv, mis
kuulub \VUgras{langustide} ja \VUgras{homaaride} \VUgras{müüjale} \textsuperscript{(102)}.

%%K t15-18-1 p
18. --- Kõik liiguvad ja teevad lärmi; kuid see ei takista selgesti kuulmast
rändava \VUgras{klaassepa} \textsuperscript{(103)} häält ega \VUgras{söekaupmehe}
\textsuperscript{(104)} hüüdu, kes on äsja mõneks hetkeks jätnud oma käru, et tarnida
\VUgras{söekott} \textsuperscript{(105)}.

\VUsaut{6.0mm}
\VUfilet{22mm}
